\documentclass[10pt, twocolumn, letterpaper]{article}

\usepackage[utf8]{inputenc}
\usepackage[T1]{fontenc}
\usepackage{mathptmx} % Times New Roman clone for standard journal feel
\usepackage{geometry}
\geometry{left=1in, right=1in, top=1in, bottom=1in, columnsep=0.25in}
\usepackage{titlesec}
\usepackage{abstract}
\usepackage{cite}
\usepackage{hyperref}

\titleformat{\section}{\normalsize\bfseries\uppercase}{\thesection.}{1em}{}
\titleformat{\subsection}{\normalsize\itshape}{\thesubsection.}{1em}{}

\title{\textbf{Operationalizing Capacity in Non-Terminal MAiD:\\The Sovereignty Model and Psychiatric Practice}}
\author{Justin Bogner\thanks{Independent Scholar. Contact: justin@example.com}}
\date{\vspace{-5ex}}

\begin{document}

\twocolumn[
  \begin{@twocolumnfalse}
    \maketitle
    \begin{abstract}
    The expansion of Medical Aid in Dying (MAiD) eligibility to include persons with primary psychiatric conditions, particularly Treatment-Resistant Depression (TRD), has provoked intense clinical debate. Opponents frequently invoke the ``depressive incapacity objection,'' arguing that the hopelessness endemic to severe depression inherently compromises decisional capacity, thereby rendering informed refusal invalid. This paper challenges that assumption. Drawing on the MacArthur Competence Assessment Tool for Treatment (MacCAT-T) and recent empirical literature, this article argues that hopelessness does not constitutively override the cognitive capacity required to evaluate medical alternatives. The paper proposes a ``Sovereignty Model'' of clinical assessment that replaces paternalistic custodial retention with a robust, two-stage capacity evaluation. Furthermore, it introduces a Structural Accountability Mechanism to address the confounding variable of socioeconomic deprivation, ensuring that clinical psychiatry functions as a facilitator of authentic autonomy rather than an administrative gatekeeper for systemic policy failures.
    \vspace{2em}
    \end{abstract}
  \end{@twocolumnfalse}
]

\section{Introduction}

The integration of Medical Aid in Dying (MAiD) into mainstream clinical practice represents a paradigm shift in the application of medical ethics. In jurisdictions evaluating or implementing non-terminal eligibility pathways---such as Canada's Track 2 framework---the psychiatric profession finds itself at the center of a profound structural conflict. On one side is the foundational medical doctrine of informed refusal: the right of a capacitated adult to decline life-sustaining treatment regardless of the physician's prognosis \cite{beauchamp2019}. On the other side is the historical orientation of psychiatry toward suicide prevention, an orientation heavily defined by institutional risk management and biological continuation \cite{jobes2016}.

When a patient with Treatment-Resistant Depression (TRD) requests MAiD, these two imperatives collide. The dominant clinical response has relied on the ``Protectionist Model,'' which utilizes suffering legibility and exhaustion of treatments as primary gatekeeping mechanisms. 

This paper argues that the Protectionist Model is clinically incoherent. In its place, it proposes the ``Sovereignty Model,'' establishing decisional capacity as the sole eligibility criterion.

\section{The Depressive Incapacity Objection}

The most formidable barrier to the Sovereignty Model is the ``depressive incapacity objection.'' This objection posits that persons with TRD cannot satisfy the \textit{appreciation} criterion of decisional capacity because hopelessness functions as a cognitive filter. While a patient might formally understand the alternatives (e.g., a novel pharmacological trial like esketamine, or ECT), their depression renders these alternatives subjectively inaccessible \cite{sullivan1994}.

\subsection{Differentiating Capacity from Optimism}

The objection conflates clinical capacity with emotional optimism. The appreciation criterion of the widely validated MacArthur Competence Assessment Tool for Treatment (MacCAT-T) requires that a patient understand how medical information applies to their specific situation \cite{appelbaum1995}. 

Empirical studies using the MacCAT-T consistently demonstrate that many patients with severe, treatment-resistant depression retain the cognitive capacity to understand, reason, and express a choice about their treatment options, even when burdened by profound hopelessness \cite{hindmarch2013, vollmann2003}. A patient who accurately states that selective serotonin reuptake inhibitors (SSRIs) have a documented response rate, that multiple adequate trials have failed, and that they decline further trials because the statistical probability of benefit does not justify the psychological cost, perfectly satisfies the appreciation criterion.

To require that the patient also \textit{feel hope} that the alternatives might succeed is to import a positive affective requirement with no basis in the Appelbaum-Grisso framework. It creates a circular standard where the primary symptom of the disease---hopelessness---becomes a permanent legal disqualifier for bodily autonomy \cite{donnelly2024}.

\section{The Sovereignty Protocol}

To operationalize autonomy safely, the Sovereignty Model replaces subjective assessments of suffering with a rigorous, two-stage capacity protocol.

\subsection{Stage 1: Decisional Capacity}

Stage 1 relies on dual independent evaluations utilizing the four elements of the MacCAT-T:
1. \textbf{Understanding:} The applicant understands the permanence of death and the realistic outcomes of all presented alternatives.
2. \textbf{Appreciation:} The applicant correctly applies the statistical and clinical realities to their own medical history.
3. \textbf{Reasoning:} The applicant can articulate a coherent account linking their situation, values, and decision.
4. \textbf{Expression:} The preference is stable across multiple independent evaluations and throughout a mandatory 90-day assessment period.

Crucially, the clinical interview tests \textit{comprehension} of alternatives, not \textit{compliance}. Informed refusal of a novel therapeutic intervention (e.g., rTMS) is as legally and clinically valid as informed refusal of an antibiotic.

\subsection{Stage 2: Process Comprehension}

Stage 2 evaluates the applicant's metacognitive awareness. It requires the applicant to demonstrate an understanding of why the procedural safeguards exist. A patient experiencing an acute, transient crisis will rarely sustain engagement with a multi-month bureaucratic procedure. The administrative friction of the 90-day wait serves as a diagnostic filter, effectively distinguishing chronic, stable ideation from acute crisis states.

\section{The Structural Accountability Mechanism}

A valid clinical concern regarding psychiatric MAiD is the risk of interpersonal or socioeconomic coercion. If a patient requests MAiD because they lack adequate housing or cannot afford disability care, authorizing the request without addressing the systemic failure weaponizes medical autonomy against the marginalized.

The Sovereignty Model institutes a Structural Accountability Mechanism. During the capacity assessment, an unaffiliated social worker must conduct an independent coercion and resource screening. If socioeconomic deprivation is identified as a primary driver, this data is mandatorily recorded and reported to the state oversight body.

This mechanism does not paternalistically deny MAiD to those in poverty---which would trap them in the very suffering they wish to escape---but instead forces the state to publicly document the structural conditions driving its citizens to choose death. It shifts the psychiatrist's role from gatekeeping systemic failures to surfacing them for public accountability.

\section{Redefining the Therapeutic Alliance}

The most profound consequence of the Sovereignty Model is the realignment of the therapeutic relationship. When a patient possesses an enforceable right to exit, the clinician is stripped of the coercive power of custodial retention.

This constraint is the precondition for true clinical effectiveness. The therapeutic relationship transitions from custodial management to a restorative alliance. The clinician's role is no longer to enforce biological continuity through liability theater (e.g., no-suicide contracts), but to actively make the patient's life subjectively worth continuing.

\section{Conclusion}

The assessment of decisional capacity in the psychiatric population is complex, but it is not impossible. The Sovereignty Model recognizes that the challenge of complex cases is an argument for rigorous, multi-evaluator assessments, not for the categorical exclusion of psychiatric patients from the fundamental right to bodily autonomy. 

By grounding MAiD eligibility in capacity rather than suffering, psychiatry can honor the informed refusal doctrine, protect vulnerable populations through structural accountability, and return to its core mission: alleviating distress through collaborative, non-coercive care.

\begin{thebibliography}{9}
\small
\bibitem{appelbaum1995} Appelbaum, P.S., \& Grisso, T. (1995). The MacArthur Treatment Competence Study. I: Mental illness and competence to consent to treatment. \textit{Law and Human Behavior}, 19(2), 105--126.
\bibitem{beauchamp2019} Beauchamp, T.L., \& Childress, J.F. (2019). \textit{Principles of biomedical ethics} (8th ed.). Oxford University Press.
\bibitem{donnelly2024} Donnelly, M., \& Campbell, A. (2024). Decisional capacity and the limits of paternalism in assisted dying law. \textit{Journal of Medical Ethics}, 50(3), 182--191.
\bibitem{hindmarch2013} Hindmarch, T., Hotopf, M., \& Owen, G. S. (2013). Depression and decision-making capacity for treatment or research: a systematic review. \textit{BMC Medical Ethics}, 14(1), 54.
\bibitem{jobes2016} Jobes, D.A. (2016). \textit{Managing suicidal risk: A collaborative approach} (2nd ed.). Guilford Press.
\bibitem{sullivan1994} Sullivan, M.D., \& Youngner, S.J. (1994). Depression, competence, and the right to refuse lifesaving medical treatment. \textit{American Journal of Psychiatry}, 151(7), 971--978.
\bibitem{vollmann2003} Vollmann, J., et al. (2003). Competence of mentally ill patients: a comparative empirical study. \textit{Psychological Medicine}, 33(8), 1463--1471.
\end{thebibliography}

\end{document}
